\chapter{MoveIt 2}
\label{ch:moveit}

\section{Launching}

\begin{lstlisting}
# move_group against the simulated controllers
ros2 launch baxter_moveit_config sim_moveit.launch.py headless:=true

# and the MotionPlanning RViz panel
ros2 launch baxter_moveit_config sim_moveit_rviz.launch.py headless:=false
\end{lstlisting}

Wait for \texttt{You can start planning now!} and a log line naming pipeline
\texttt{ompl}. The launch is readiness-gated: it waits for the simulation to be
up before starting \texttt{move\_group}, because a \texttt{move\_group} that
loads before the controllers exist comes up without them and then fails to
execute anything it plans.

\begin{figure}[h]
\centering
\graphicsorplaceholder{figures/moveit_rviz.png}{0.92\linewidth}{the RViz
MotionPlanning panel with the \texttt{both\_arms} group selected}
\caption{The MotionPlanning panel loaded correctly: a populated planning group,
Global Status Ok, and a 6-DOF interactive marker on each gripper. Note the
Velocity Scaling of 0.30 --- that is the RViz default and is fine in simulation.
On hardware it is too fast; see Section~\ref{sec:moveit-hardware}.}
\label{fig:moveitrviz}
\end{figure}

\section{Groups}

Five groups are defined: \texttt{left\_arm}, \texttt{right\_arm},
\texttt{both\_arms}, and a hand group per side. Named states
\texttt{left\_neutral} and \texttt{right\_neutral} match the simulation's initial
joint values exactly --- CI asserts that equality, because a neutral pose that
drifts from the model's start state makes every "return to neutral" a silent
small motion.

Planning uses OMPL with \texttt{RRTConnectkConfigDefault} as the default for all
three arm groups.

\section{The example clients}

\begin{lstlisting}
ros2 run baxter_examples moveit_left_tiny
ros2 run baxter_examples moveit_tiny --ros-args -p group:=right_arm
ros2 run baxter_examples moveit_tiny --ros-args -p group:=both_arms
\end{lstlisting}

Same contract as the direct client: fresh state, bounded reversible target,
verified against a fresh reading, returned to start. Add
\texttt{-p cancel\_after\_sec:=1.0} to exercise cancellation.

\begin{notebox}[These clients refuse to touch hardware]
\texttt{moveit\_left\_tiny} checks that \texttt{/move\_group} is running on
simulated time and refuses to send a goal otherwise. It is deliberate: a sim
example that could be pointed at a real robot by changing one launch file is a
sim example waiting to move an arm nobody was watching. Hardware MoveIt execution
is driven from the RViz panel by a person, as described in
Section~\ref{sec:moveit-hardware}.
\end{notebox}

\section{Cartesian goals and IK}

\begin{lstlisting}
# absolute point
ros2 run baxter_examples moveit_pose --ros-args \
    -p group:=left_arm -p x:=0.55 -p y:=0.25 -p z:=0.15

# relative nudge
ros2 run baxter_examples moveit_pose --ros-args \
    -p group:=left_arm -p delta_z:=0.05

# plan only, do not execute
ros2 run baxter_examples moveit_pose --ros-args \
    -p group:=left_arm -p delta_z:=0.05 -p plan_only:=true
\end{lstlisting}

Targets are position-only against a 3\,cm tolerance sphere in \texttt{torso} by
default, so gripper orientation is whatever IK picks. Absolute and relative forms
are mutually exclusive. \texttt{plan\_only} is the form to use first whenever
\texttt{move\_group} is attached to anything that can move.

For pure IK queries without motion:

\begin{lstlisting}
ros2 run baxter_examples ik_service_client --ros-args -p limb:=left
ros2 run baxter_examples ik_service_client --ros-args -p limb:=left -p x:=9.0
\end{lstlisting}

The second is deliberately unreachable and must exit non-zero with
\texttt{INVALID POSE}. A solver that never reports failure is not a solver, and
the verification script asserts that failure explicitly.

\section{The collision matrix}

\texttt{config/baxter.srdf} carries exactly 54 \texttt{disable\_collisions} pairs.
That is sparse for a robot this size, and two of those pairs were added by hand:

\begin{center}
\texttt{left\_upper\_elbow} $\leftrightarrow$ \texttt{left\_upper\_shoulder}\\
\texttt{right\_upper\_elbow} $\leftrightarrow$ \texttt{right\_upper\_shoulder}
\end{center}

Those links are two apart in the chain and graze by about 1.7\,mm in the pose the
untuck tool leaves the robot in. Without the pairs disabled, \texttt{move\_group}
returns $-10$ \texttt{START\_STATE\_IN\_COLLISION} and refuses to plan from the
robot's own resting pose.

The MoveIt Setup Assistant cannot rediscover them. Its sampling does not converge
for this model --- runs from ten thousand to two million trials produced matrices
between 263 and 311 pairs that disabled head, gripper and screen pairs that should
not be disabled, while still missing these two. So a regenerated matrix is a
regression, and the checked-in one stays.

\begin{safetybox}[If you regenerate the SRDF]
Run \texttt{scripts/check\_srdf.py diff OLD NEW} before and after. CI asserts
both pairs are present and that the total is exactly 54; the check exists because
this failure only shows up on hardware, at the moment you most want to plan.
\end{safetybox}
